\documentclass[11pt,a4paper]{article}
\usepackage[margin=25mm]{geometry}
\usepackage{amsmath,amssymb,amsthm,hyperref}
\setlength{\emergencystretch}{2em}
\newtheorem{theorem}{Theorem}
\newcommand{\gen}[1]{\langle#1\rangle}
\newcommand{\diag}{\operatorname{diag}}
\title{Completion of the binary omission with intransitive subomissions}
\author{Human Compression project: a component of the L044 recovery}
\date{}
\begin{document}\maketitle
Let $G=A_n$ in its natural action, with distinguished generators
$\rho_0,\ldots,\rho_{r-1}$ forming a string C-marking. We acquire
the exact finite maxima below degree fifteen, the complete transitive
rank theorem, the completed pair-block theorem and the complete
contracted-forest and neighbour components. The pair-block theorem's
smaller-degree induction is stated explicitly in the result below.

\begin{theorem}
Suppose $n\ge12$, a transitive maximal omission $H=G_i$
preserves two blocks of size $m=n/2$, and every distinguished
maximal omission $H_j=G_{i,j}$ is intransitive. Assume the
alternating rank bound $\lfloor(q-1)/2\rfloor$ in every smaller
degree $12\le q<n$. Then
\[
                       r\le\lfloor(n-1)/2\rfloor.
\]
\end{theorem}
\begin{proof}
The exact finite maxima settle $n\le14$. Suppose $n\ge15$
and that the displayed bound fails, so $r\ge m$. Since $H$
is transitive, some distinguished mark $l$ interchanges its two
blocks. It is a product of $m$ transpositions, so evenness makes
$m$ even. Consequently $n\ge16$ and $m\ge8$.

Let $C$ contain the other distinguished marks of $H$ commuting
with $l$, and let $R$ contain the remaining marks. String relations
give $|R|\le2$. First suppose $R$ is empty. Choose coordinates
matching the two rows through $l$. Then
$H\le\diag(S_m)\times\gen l$, and the projection of
$\gen C$ on the $l$-pairs is faithful: its kernel is contained
in $\gen l$, while the actual intersection property gives
$\gen C\cap\gen l=1$. This image $P$ is transitive on the
pairs, because $H$ is transitive and adjoining $l$ changes no
pair index. It has the actual inherited string C-marking of rank
$r-2\ge m-2$.

For $m\ge10$, the complete transitive rank theorem forces
$P=S_m$ or $A_m$, since $m-2>m/2+1$. At $m=8$, the complete
proper transitive boundary has no rank-six exception, giving the
same conclusion. Thus the pair action is primitive and the
$l$-pair system is maximal. The completed pair-block theorem
for the parent $A_{2m}$ excludes it under the present explicit
smaller-degree induction. This establishes maximality before
that theorem is invoked; preservation of pairs alone would not
have sufficed.

Now suppose $R\ne\varnothing$. Select one actual edge for each
mark of $H$ other than $l$, joining different orbits of its actual
omission, and contract the $l$-pairs. The complete forest theorem
gives $r\le m$: a noncommuting neighbour has at least two
distinct contracted edges, which cannot both occur within a
connected selected forest. Thus the assumed failure forces
$r=m$, with exactly two selected components.

If $|R|=1$, the complete one-neighbour theorem excludes the case
where the neighbour preserves the binary rows. If it interchanges
them, that theorem changes the distinguished block swap to this
neighbour and gives exactly two noncommuting neighbours, with
the same $H$, omitted label and binary partition. All other
standing hypotheses are unchanged. Thus it remains only to treat
$|R|=2$.

Write these two neighbours as $u,v$. The complete common bridge
theorem gives rank three for their full contracted graph. The
four-pair reduction makes this graph connected on exactly four
pairs. If $u,v$ both preserve the rows, the complete preserving
neighbour theorem excludes the configuration by its actual
classification, subgroup intersections and endpoint centralizers.
If at least one interchanges the rows, the complete swapping
neighbour theorem constructs a central pair system, proves it
maximal, and applies the completed pair-block theorem. Its only
alternating induction premise is the present bound in strictly
smaller degrees. The one-neighbour reduction is used only to
enter this two-neighbour case; neither two-neighbour component
invokes a completed one-neighbour exclusion, so this routing is
not circular.

Every possibility for $R$ is excluded. Hence a rank failure is
impossible, proving the theorem.
\end{proof}

\paragraph{Attribution and scope.}
This completes the branch corresponding to source Proposition~3.5
of Cameron--Fernandes--Leemans--Mixer, conditional on the stated
smaller-degree induction. The eligible-edge strategy is borrowed;
the actual support, attachment, local permutation and centralizer
arguments are retained in full in their component manuscripts.
The empty-$R$ maximality interface and this assembly were supplied
by the continuation investigator. The distinct authorship and
review contributions to each component remain in their records.
This is an assembly of the fixed principal-source contract, not
a new conjecture claim or a whole-source disposition.
\end{document}
